\documentclass[11pt]{article}
\usepackage[margin=1.1in]{geometry}
\usepackage{amsmath,amssymb,amsthm}
\usepackage{enumitem}
\usepackage{booktabs}
\usepackage[colorlinks=true,linkcolor=blue,citecolor=blue,urlcolor=blue]{hyperref}

\newtheorem{theorem}{Theorem}[section]
\newtheorem{proposition}[theorem]{Proposition}
\newtheorem{remark}[theorem]{Remark}
\theoremstyle{definition}
\newtheorem{definition}[theorem]{Definition}

\newcommand{\pP}{\textbf{[P]}}
\newcommand{\pC}{\textbf{[C]}}
\newcommand{\pO}{\textbf{[O]}}
\newcommand{\Tr}{\operatorname{Tr}}

\title{\textbf{The New Standard Model for Unified Physics}\\[4pt]
\large The Entanglement-Certified Effective Lagrangian of Matter,
Gauge Fields, and Riemann--Cartan Gravity}
\author{Robert W.\ McGwier\\
\small CTO and Staff Scientist, Cohere Technology Group, Herndon, VA}
\date{August 7, 2026 --- Version 2}

\begin{document}
\maketitle

\begin{abstract}
\noindent
We assemble, term by term, the complete effective Lagrangian of known
physics coupled to gravity with torsion: the $SU(3)\times SU(2)\times U(1)$
Standard Model minimally coupled, through vielbein and independent spin
connection, to Einstein--Cartan--Sciama--Kibble gravity with cosmological
constant. We call this the \emph{New Standard Model} (NSM) and state
precisely the sense in which it unifies: not by embedding the gauge group
in a larger symmetry, but by certification---its gravitational sector is,
at linear order, the unique response permitted by the entanglement first
law \cite{McGwier2026a}; its torsion coupling is the unique structure
surviving the modular selection rule \cite{McGwier2026b}; its superposition
principle for geometry is a theorem rather than a postulate; and the
elimination of torsion yields a single parameter-free prediction, a
universal gravitational-strength axial--axial contact interaction among
all fermion species. We derive the full Lagrangian, prove its certified
properties, list its falsifiable consequences, and---under the standing
verification discipline---itemize with equal care everything the model
does \emph{not} explain. Every material claim is labeled \pP{}
(proved/derived), \pC{} (conjectured with support), or \pO{} (open).
\end{abstract}

\section{What ``New Standard Model'' Means Here}

The Standard Model (SM) of particle physics
\cite{Glashow1961,Weinberg1967,Salam1968} plus general relativity (GR)
constitutes the empirically confirmed description of nature. Their
juxtaposition is usually presented as provisional bookkeeping awaiting a
final unified theory. This paper takes a different, sharper position,
inherited from its two predecessors \cite{McGwier2026a,McGwier2026b}:

\begin{definition}[The New Standard Model]
The NSM is the effective field theory defined by the action
\eqref{eq:NSM} below---the full SM coupled to Einstein--Cartan gravity in
first-order variables---\emph{together with} the certification that its
gravitational dynamics at linear order is equivalent to the entanglement
first law of the underlying quantum state, and the selection rule that its
torsion structure is the unique one consistent with that law.
\end{definition}

``Unified'' is therefore used in a precise and limited sense: gravity,
inertia, torsion, and the superposition principle for geometry are not
independent postulates adjoined to quantum matter; within the certified
regime they are consequences of the entanglement structure of the matter
state \cite{Jacobson1995,Faulkner2014,McGwier2026a}. What is \emph{not}
claimed---gauge unification, mass origins, a final theory---is itemized in
Section~7 with the same care as the claims. The framework's own logic
forbids finality: the Lagrangian below is emergent bookkeeping for a more
fundamental description, and the position defended throughout this program
is that every such description admits a deeper layer.

\section{Field Content and Variables}

Gravitational variables: coframe $e^a=e^a{}_\mu dx^\mu$ and independent
Lorentz connection $\omega^{ab}{}_\mu$, with torsion
$T^a=de^a+\omega^a{}_b\wedge e^b$ and curvature
$R^{ab}=d\omega^{ab}+\omega^a{}_c\wedge\omega^{cb}$.

Matter content: the confirmed SM fields---gauge potentials
$G^A_\mu, W^I_\mu, B_\mu$ of $SU(3)_c\times SU(2)_L\times U(1)_Y$; three
generations of chiral fermions $\{Q_L,u_R,d_R,L_L,e_R\}$ in the standard
representations; the Higgs doublet $H$ with its measured quartic potential
and 125\,GeV excitation \cite{ATLAS2012,CMS2012}; and neutrino mass terms
required by oscillation data \cite{SuperK1998}, implemented either by
right-handed singlets $\nu_R$ (Dirac) or by the dimension-five Weinberg
operator (Majorana)---the choice is empirically open. \pP{} for the
content; Dirac-vs-Majorana \pO{}.

Fermions couple to gravity through the vielbein and the full connection:
\begin{equation}
D_\mu\psi=\Big(\partial_\mu+\tfrac{1}{4}\omega_{ab\mu}\gamma^{ab}
+ig_s G_\mu+ig W_\mu+ig' Y B_\mu\Big)\psi\,.
\end{equation}
Gauge field strengths are built with the exterior derivative alone,
$F=dA+A\wedge A$: minimal substitution of the torsionful connection into
$F$ would violate gauge invariance, so gauge fields do not couple to
torsion \cite{HehlRMP1976,Shapiro2002}. Scalars carry no spin and likewise
do not source torsion. \pP{} Hence \emph{only fermions talk to torsion}, a
structural fact with consequences below.

\section{The Lagrangian}

\begin{equation}
\boxed{\;
S_{\mathrm{NSM}}
=\frac{1}{4\kappa}\!\int\!\epsilon_{abcd}\Big(e^a\wedge e^b\wedge R^{cd}
-\frac{\Lambda}{6}\,e^a\wedge e^b\wedge e^c\wedge e^d\Big)
\;+\;\int d^4x\;e\,\mathcal{L}_{\mathrm{SM}}[e,\omega,\psi,A,H]\;}
\label{eq:NSM}
\end{equation}
with $\kappa=8\pi G$, $e=\det e^a{}_\mu$, and
\begin{align}
\mathcal{L}_{\mathrm{SM}}
=\;&-\tfrac{1}{4}G^A_{\mu\nu}G^{A\mu\nu}
   -\tfrac{1}{4}W^I_{\mu\nu}W^{I\mu\nu}
   -\tfrac{1}{4}B_{\mu\nu}B^{\mu\nu}
\;+\;\sum_f \tfrac{i}{2}\,\bar\psi_f\gamma^a e^\mu{}_a
      \overleftrightarrow{D}_\mu\psi_f
\nonumber\\[2pt]
&+\,|D_\mu H|^2 - V(H)
\;-\;\Big(\bar Q_L Y_u \tilde H u_R + \bar Q_L Y_d H d_R
      + \bar L_L Y_e H e_R + \mathrm{h.c.}\Big)
\;+\;\mathcal{L}_\nu\,,
\label{eq:LSM}
\end{align}
where $\mathcal{L}_\nu$ is the Dirac or Weinberg neutrino term. The purely
gravitational sector of \eqref{eq:NSM} reproduces
$\tfrac{1}{2\kappa}\!\int\!\sqrt{-g}\,(R-2\Lambda)$ when torsion vanishes.
\pP{}

\subsection{Field equations}

Independent variation \cite{Kibble1961,Sciama1964,HehlRMP1976}:
\begin{align}
\delta e:&\quad
\tfrac{1}{2}\epsilon_{abcd}\,e^b\wedge\!\Big(R^{cd}
-\tfrac{\Lambda}{3}e^c\wedge e^d\Big)=\kappa\,\tau_a\,,
\label{eq:einsteinsector}\\
\delta\omega:&\quad
\epsilon_{abcd}\,e^a\wedge T^b=\kappa\,\mathcal{S}_{cd}\,,
\qquad
\mathcal{S}\ \text{sourced by fermions only.}
\label{eq:cartansector}
\end{align}
The Cartan equation \eqref{eq:cartansector} is algebraic; torsion is
locked pointwise to the total fermionic spin density and vanishes in
vacuum. \pP{}

\subsection{Exact elimination of torsion: the one new interaction}

Each SM fermion's canonical spin current is totally antisymmetric and dual
to its axial current, $s_f^{\lambda\mu\nu}=
-\tfrac14\epsilon^{\lambda\mu\nu\rho}J_{5\rho}^{(f)}$
\cite{HehlDatta1971,McGwier2026b}. Solving \eqref{eq:cartansector} and
substituting back yields ordinary Riemannian gravity plus a single
contact term \cite{HehlDatta1971}:
\begin{equation}
\boxed{\;
\mathcal{L}_{\mathrm{HD}}
=-\frac{3\kappa}{16}\,
\mathcal{J}_5^{\mu}\,\mathcal{J}_{5\mu}\,,
\qquad
\mathcal{J}_5^\mu\equiv\sum_{f\,\in\,\mathrm{all\ fermions}}
\bar\psi_f\gamma^5\gamma^\mu\psi_f\;}
\label{eq:HD}
\end{equation}
---a \emph{universal, flavor-blind, gravitational-strength axial--axial
interaction of the total axial current with itself}, including
cross-species terms (quark--lepton, quark--neutrino) with no free
parameter: the coefficient is fixed by $G$. This is the only interaction
in the NSM beyond the SM and Einstein gravity, and it is an output, not an
input. \pP{}

\section{The Certifications: What Makes This ``New''}

\begin{theorem}[Gravity sector certified by entanglement]
At linear order about the vacuum, the metric dynamics of \eqref{eq:NSM} is
equivalent to the first law of entanglement
$\delta S_B=\delta\langle H_B\rangle$ imposed on all ball-shaped regions:
curvature is the unique consistent response to perturbations of vacuum
entanglement, with inertia and the Newtonian limit contained in the same
statement. \pP{} within the holographic regime
\cite{Faulkner2014,Lashkari2014,McGwier2026a}; thermodynamic derivations
extend the statement to the full nonlinear torsionful equations under
Clausius assumptions \cite{Jacobson1995,DLP2017,JacobsonMohd2015}.
\end{theorem}

\begin{theorem}[Torsion structure certified by modular selection]
The torsion coupling of \eqref{eq:NSM} is minimal and axial. This is not a
simplicity choice: spin-current sources with mixed-symmetry structure,
coupling through non-minimal matter, are relevant deformations that
destroy the conformal/modular structure on which the first law rests; the
entanglement-consistent sector is exactly the conserved-current (axial,
minimal) sector realized in \eqref{eq:NSM}--\eqref{eq:HD}. \pP{} at the
level of the dimensional and symmetry argument of \cite{McGwier2026b}.
\end{theorem}

\begin{theorem}[Superposition built in]
The NSM requires no quantization of the metric as an independent field and
no collapse mechanism: weak-field superposition of gravitational fields is
a corollary of the linearity of the first law, and superpositions of
macroscopically distinct geometries are superpositions of area-operator
eigenstates in the sense of holographic error correction. \pP{} within
AdS/CFT \cite{Harlow2017,McGwier2026a}; beyond it \pO{}.
\end{theorem}

\begin{remark}[Status of the matter sector]
No analogous certification exists for the gauge and Yukawa structure: the
SM content of \eqref{eq:LSM} enters as empirical input. Whether gauge
dynamics, chirality, and three-generation structure admit an
entanglement-theoretic derivation is the deepest open question this
program can pose. \pO{}
\end{remark}

\section{Falsifiable Consequences}

\begin{enumerate}[itemsep=2pt]
\item \textbf{Universal axial contact term.} Eq.~\eqref{eq:HD} predicts
axial--axial four-fermion interactions of strength $\sim G$, i.e.\
suppressed by $M_{\mathrm{Pl}}^{-2}$: absolutely fixed, absolutely tiny,
and far below collider contact-interaction sensitivity---but nonzero, and
cross-species. Laboratory and astrophysical bounds on background torsion
and torsion-induced Lorentz-violating couplings exist and are consistent
\cite{KRT2008,Shapiro2002}. \pP{} for the prediction; detection prospects
at accessible energies: none. \pP{}
\item \textbf{High-density bounce.} At spin densities where
$\kappa\,\mathcal{J}_5^2$ rivals the energy density, \eqref{eq:HD} is
repulsive and can replace the initial singularity with a bounce
\cite{Trautman1973,Poplawski2010}. \pC{}
\item \textbf{Gravitationally mediated entanglement.} The NSM requires
that gravity carry quantum coherence; the Bose et al./Marletto--Vedral
witness experiments \cite{Bose2017,MarlettoVedral2017} are a
discriminator against collapse alternatives. A confirmed null result at
sufficient sensitivity falsifies the framework's gravitational sector.
\pP{} for the logical structure; \pC{} pending data.
\item \textbf{Infrared deviations.} If the de Sitter extension of the
entanglement derivation follows Verlinde's emergent-gravity route, the
NSM's gravitational law acquires infrared modifications at the Milgrom
scale without particle dark matter \cite{Verlinde2017,Milgrom1983}. \pC{};
this is the sector where the program meets galactic-dynamics data.
\end{enumerate}

\section{Relation to the Predecessor Documents}

Paper I \cite{McGwier2026a} certified curvature from entanglement and
derived the Einstein--Cartan extension; Paper II \cite{McGwier2026b}
resolved the torsionful first law into its theorem and counterexample
sectors, producing the selection rule used in Theorem 4.2. The present
paper is the assembly: the unique Lagrangian those results select when the
empirical matter content is installed. The three documents form one
program: \emph{certify, select, assemble}.

\section{What This Model Does Not Do}

Stated with the same rigor as the claims. The NSM does \emph{not}:
(i) unify the gauge couplings or predict their values;
(ii) explain fermion masses, mixings, generations, or the hierarchy
problem; (iii) resolve strong CP; (iv) supply a dark-matter particle---it
offers instead the conjectural emergent-gravity alternative of item 5.4;
(v) determine the cosmological constant, though the framework classifies
the usual fine-tuning argument as an artifact of Wilsonian assumptions
\pC{}; (vi) provide its own ultraviolet completion---by construction: the
NSM is asserted to be \emph{emergent} bookkeeping for a more fundamental
description whose full form is \pO{}, with the certified regime anchored
in holography. Finality is not claimed and, on this program's own
epistemology, should not be claimable. All \pO{}/\pC{} as marked.

\section{Certified Summary}

The New Standard Model is the action \eqref{eq:NSM}: confirmed matter,
minimally coupled to first-order gravity with cosmological constant, its
torsion eliminable in favor of the parameter-free universal axial
interaction \eqref{eq:HD}. Its gravitational dynamics is entanglement
bookkeeping \pP{}; its torsion structure is modular-selected \pP{}; its
geometry superposes as a theorem \pP{}; its matter content is empirical
input \pP{}; its unexplained structures are enumerated, not obscured. It
is offered not as a final theory but as the current certified layer---the
next turtle, load-bearing and inspected.

\section*{Acknowledgments}
The author acknowledges the use of Claude (Anthropic, claude.ai) for
assistance in drafting, literature verification, and \LaTeX{} preparation.
All technical claims and epistemic classifications were reviewed and
verified by the author, who bears sole responsibility for the content.

\begin{thebibliography}{99}\small

\bibitem{Glashow1961}
S.~L.~Glashow, ``Partial symmetries of weak interactions,''
Nucl.\ Phys.\ \textbf{22}, 579 (1961).

\bibitem{Kibble1961}
T.~W.~B.~Kibble, ``Lorentz invariance and the gravitational field,''
J.\ Math.\ Phys.\ \textbf{2}, 212 (1961).

\bibitem{Sciama1964}
D.~W.~Sciama, ``The physical structure of general relativity,''
Rev.\ Mod.\ Phys.\ \textbf{36}, 463 (1964).

\bibitem{Weinberg1967}
S.~Weinberg, ``A model of leptons,''
Phys.\ Rev.\ Lett.\ \textbf{19}, 1264 (1967).

\bibitem{Salam1968}
A.~Salam, ``Weak and electromagnetic interactions,'' in
\emph{Elementary Particle Theory}, ed.\ N.~Svartholm
(Almqvist \& Wiksell, Stockholm, 1968), p.~367.

\bibitem{HehlDatta1971}
F.~W.~Hehl and B.~K.~Datta, ``Nonlinear spinor equation and asymmetric
connection in general relativity,''
J.\ Math.\ Phys.\ \textbf{12}, 1334 (1971).

\bibitem{Trautman1973}
A.~Trautman, ``Spin and torsion may avert gravitational singularities,''
Nature Phys.\ Sci.\ \textbf{242}, 7 (1973).

\bibitem{HehlRMP1976}
F.~W.~Hehl, P.~von der Heyde, G.~D.~Kerlick and J.~M.~Nester, ``General
relativity with spin and torsion,''
Rev.\ Mod.\ Phys.\ \textbf{48}, 393 (1976).

\bibitem{Milgrom1983}
M.~Milgrom, ``A modification of the Newtonian dynamics,''
Astrophys.\ J.\ \textbf{270}, 365 (1983).

\bibitem{Jacobson1995}
T.~Jacobson, ``Thermodynamics of spacetime: the Einstein equation of
state,'' Phys.\ Rev.\ Lett.\ \textbf{75}, 1260 (1995).

\bibitem{SuperK1998}
Super-Kamiokande Collaboration (Y.~Fukuda et al.), ``Evidence for
oscillation of atmospheric neutrinos,''
Phys.\ Rev.\ Lett.\ \textbf{81}, 1562 (1998).

\bibitem{Shapiro2002}
I.~L.~Shapiro, ``Physical aspects of the space-time torsion,''
Phys.\ Rept.\ \textbf{357}, 113 (2002), arXiv:hep-th/0103093.

\bibitem{KRT2008}
V.~A.~Kosteleck\'y, N.~Russell and J.~D.~Tasson, ``Constraints on torsion
from bounds on Lorentz violation,''
Phys.\ Rev.\ Lett.\ \textbf{100}, 111102 (2008), arXiv:0712.4393.

\bibitem{Poplawski2010}
N.~J.~Pop\l awski, ``Cosmology with torsion,''
Phys.\ Lett.\ B \textbf{694}, 181 (2010), arXiv:1007.0587.

\bibitem{ATLAS2012}
ATLAS Collaboration, ``Observation of a new particle in the search for
the Standard Model Higgs boson,''
Phys.\ Lett.\ B \textbf{716}, 1 (2012), arXiv:1207.7214.

\bibitem{CMS2012}
CMS Collaboration, ``Observation of a new boson at a mass of 125 GeV,''
Phys.\ Lett.\ B \textbf{716}, 30 (2012), arXiv:1207.7235.

\bibitem{FLM2013}
T.~Faulkner, A.~Lewkowycz and J.~Maldacena, ``Quantum corrections to
holographic entanglement entropy,'' JHEP \textbf{11}, 074 (2013).

\bibitem{Lashkari2014}
N.~Lashkari, M.~B.~McDermott and M.~Van Raamsdonk, ``Gravitational
dynamics from entanglement `thermodynamics','' JHEP \textbf{04}, 195
(2014), arXiv:1308.3716.

\bibitem{Faulkner2014}
T.~Faulkner, M.~Guica, T.~Hartman, R.~C.~Myers and M.~Van Raamsdonk,
``Gravitation from entanglement in holographic CFTs,''
JHEP \textbf{03}, 051 (2014), arXiv:1312.7856.

\bibitem{JacobsonMohd2015}
T.~Jacobson and A.~Mohd, ``Black hole entropy and Lorentz-diffeomorphism
Noether charge,'' Phys.\ Rev.\ D \textbf{92}, 124010 (2015).

\bibitem{Bose2017}
S.~Bose et al., ``Spin entanglement witness for quantum gravity,''
Phys.\ Rev.\ Lett.\ \textbf{119}, 240401 (2017).

\bibitem{MarlettoVedral2017}
C.~Marletto and V.~Vedral, ``Gravitationally induced entanglement between
two massive particles is sufficient evidence of quantum effects in
gravity,'' Phys.\ Rev.\ Lett.\ \textbf{119}, 240402 (2017).

\bibitem{Harlow2017}
D.~Harlow, ``The Ryu--Takayanagi formula from quantum error correction,''
Commun.\ Math.\ Phys.\ \textbf{354}, 865 (2017).

\bibitem{Verlinde2017}
E.~Verlinde, ``Emergent gravity and the dark universe,''
SciPost Phys.\ \textbf{2}, 016 (2017), arXiv:1611.02269.

\bibitem{DLP2017}
R.~Dey, S.~Liberati and D.~Pranzetti, ``Spacetime thermodynamics in the
presence of torsion,'' Phys.\ Rev.\ D \textbf{96}, 124032 (2017).

\bibitem{McGwier2026a}
R.~W.~McGwier, ``Entanglement as the origin of spacetime curvature: a
verified synthesis of the first-law derivation of the Einstein
equations,'' Version 3, Zenodo (2026), DOI:
\href{https://doi.org/10.5281/zenodo.21821890}{10.5281/zenodo.21821890}.

\bibitem{McGwier2026b}
R.~W.~McGwier, ``Spin-current sources and ball modular Hamiltonians:
resolving the torsionful first law,''
Zenodo (2026), DOI:
\href{https://doi.org/10.5281/zenodo.21824934}{10.5281/zenodo.21824934}.

\bibitem{McGwier2026d}
R.~W.~McGwier, ``Condition NL for finite balls: a spectral reduction, the unconditional metric sector, and a universal kinematic residue,''
Zenodo (2026), DOI:
\href{https://doi.org/10.5281/zenodo.21836572}{10.5281/zenodo.21836572}.

\bibitem{McGwier2026e}
R.~W.~McGwier, ``The Hehl--Datta coefficient from relative-entropy positivity,''
Zenodo (2026), DOI:
\href{https://doi.org/10.5281/zenodo.21836757}{10.5281/zenodo.21836757}.

\bibitem{McGwier2026f}
R.~W.~McGwier, ``Emergent gravity and a new Lagrangian: a theory of everything we know,''
Zenodo (2026), DOI to be assigned.

\end{thebibliography}

\end{document}
